\documentclass[10pt]{article}
\usepackage[margin=0.95in]{geometry}
\usepackage{amsmath,amssymb,amsthm}
\usepackage[T1]{fontenc}
\usepackage{lmodern}
\usepackage{microtype}
\usepackage[hidelinks]{hyperref}
\newtheorem{theorem}{Theorem}
\newtheorem{corollary}[theorem]{Corollary}
\newcommand{\Tr}{\operatorname{Tr}}
\newcommand{\diag}{\operatorname{diag}}
\newcommand{\MS}{\mathcal M}
\title{A counterexample to Bu--Gu--Jaffe's expected\protect\\ mixed-state inverse sumset estimate}
\author{Haobo Ma\thanks{ChronoAI Pte Ltd, Singapore.}\and
Wenlin Zhang\thanks{National University of Singapore; ChronoAI Pte Ltd, Singapore.}}
\date{1 November 2026}
\begin{document}
\maketitle
\begin{abstract}
We refute the mixed-state inverse sumset estimate expected in Remark~45,
equation~(58), of Bu, Gu and Jaffe, \emph{Quantum Ruzsa Divergence to Quantify
Magic}~\cite[arXiv:2401.14385v3]{BGJ}.
A computationally diagonal equal mixture of two basis states in dimension seven has
quantum-doubling constant $\sqrt2$ and relative entropy $\log(7/2)$ from its
mean state, while the proposed bound is at most $\tfrac12\log7$.
The contradiction uses exact inequalities and persists for every tensor power.
For the same family, we verify that the weaker estimate~(59) has a positive
denominator and holds strictly, and give the equality case of the companion
entropy-growth statement for computationally diagonal states.
\end{abstract}

\section{The expected estimate and its conventions}

Bu, Gu and Jaffe extend entropic inverse sumset theory to discrete quantum
convolution~\cite{BGJ}. Their Remark~45 expects a mixed-state extension of
the pure-state bound in Theorem~44. We show that its displayed estimate~(58)
fails already for a diagonal state, so the obstruction is classical rather
than a consequence of noncommutativity.

Let $d$ be an odd prime and $s,t\in\mathbb F_d\setminus\{0\}$ satisfy
$s^2+t^2=1$. On computational basis vectors, the unitary and convolution are
\begin{equation}\label{eq:conv}
 U_{s,t}|i,j\rangle=|si+tj,-ti+sj\rangle,\qquad
 \rho\boxtimes_{s,t}\sigma=\Tr_2\bigl[U_{s,t}(\rho\otimes\sigma)U_{s,t}^*\bigr].
\end{equation}
For $n$ qudits, this map acts coordinatewise. All logarithms are natural;
$S(\rho)=-\Tr(\rho\log\rho)$, with $0\log0=0$.
We use $D(\rho\Vert\sigma)=\Tr\rho(\log\rho-\log\sigma)$ when
$\operatorname{supp}\rho\subseteq\operatorname{supp}\sigma$, and $+\infty$
otherwise.
We use the source's characteristic function
$\Xi_\rho(p,q)=\Tr(\rho w(-p,-q))$, where
$Z|j\rangle=\omega^j|j\rangle$, $X|j\rangle=|j+1\rangle$,
$\omega=e^{2\pi i/d}$ and
$w(p,q)=\omega^{-pq/2}Z^pX^q$; division by two in this exponent means
multiplication by $2^{-1}$ in $\mathbb F_d$.
The mean state $\MS(\rho)$ retains precisely the characteristic coefficients
of modulus one~\cite[Definitions~6--8 and 10]{BGJ}. Write
\begin{align*}
 R(\rho)&=\operatorname{rank}\MS(\rho),&
 \lambda&=(1-\operatorname{MG}(\rho))^2,\\
 \delta_q[\rho]&=\exp\bigl(S(\rho\boxtimes\rho)-S(\rho)\bigr),&
 \operatorname{MG}(\rho)&=1-\max_{0<|\Xi_\rho(x)|<1}|\Xi_\rho(x)|,
\end{align*}
with the source's convention $\operatorname{MG}=0$ when the maximizing set is
empty. Its expected estimate, when $1<\delta_q[\rho]\le C$, is
\begin{equation}\label{eq:expected}
 D(\rho\Vert\MS(\rho))\ \le\
 \frac{\log R(\rho)}{\log R(\rho)-\log\bigl[1+\lambda
 (R(\rho)\Tr\rho^2-1)\bigr]}\log C.
 \tag{BGJ 58}
\end{equation}
Our examples satisfy strict positivity of every scalar logarithm argument in
the displayed estimates and of their denominators, as well as the source's
zero-mean condition: $\Xi_{\MS(\rho)}$ takes values in $\{0,1\}$.

The mean-state framework was developed in~\cite{PNAS,DQG}; the classical entropic
inverse problem appears in Tao~\cite{Tao}. Recent work of Goodarzi develops
submodular entropy bounds for compatible quantum convolutions~\cite[Theorem~4.1]{Goodarzi}.
The question here is the specific quantitative mixed-state estimate
\eqref{eq:expected}. The single-qudit witness below is the construction in the
trureturing source~\cite{Source}. We give its elementary exposition and derive
the tensor-product and equation-(59) consequences below.

\section{An exact counterexample}

\begin{theorem}\label{thm:counterexample}
Take $d=7$, $s=t=2$ and
$\rho=\tfrac12(|0\rangle\langle0|+|1\rangle\langle1|)$.
Then \eqref{eq:expected} is false for $C=\sqrt2$.
In fact, its right-hand side is at most $\tfrac12\log7$ even if the actual
$\lambda$ is replaced by any number in $[0,1]$, whereas
$D(\rho\Vert\MS(\rho))=\log(7/2)>\tfrac12\log7$.
\end{theorem}

\begin{proof}
The parameters are admissible since $2^2+2^2=8=1$ in $\mathbb F_7$.
For a diagonal state, \eqref{eq:conv} is the law of $2I+2J$ modulo seven,
where $I,J$ are independent and uniform on $\{0,1\}$. Thus
\begin{equation}\label{eq:output}
 \rho\boxtimes\rho
 =\tfrac14|0\rangle\langle0|+\tfrac12|2\rangle\langle2|
   +\tfrac14|4\rangle\langle4|.
\end{equation}
Consequently
\begin{equation}\label{eq:entropies}
 S(\rho)=\log2,\qquad S(\rho\boxtimes\rho)=\tfrac32\log2,
 \qquad \delta_q[\rho]=\sqrt2>1,\qquad \Tr\rho^2=\tfrac12.
\end{equation}
If $q\ne0$, $w(-p,-q)$ has zero diagonal and hence
$\Xi_\rho(p,q)=0$. On the clock axis,
\begin{equation}\label{eq:fourier}
 \Xi_\rho(p,0)=\frac{1+\omega^{-p}}2,\qquad
 |\Xi_\rho(p,0)|=\left|\cos\frac{\pi p}{7}\right|.
\end{equation}
The modulus is one only for $p=0$. Characteristic-function inversion therefore
gives
\begin{equation}\label{eq:mean}
 \MS(\rho)=I_7/7,\qquad R(\rho)=7,\qquad
 D(\rho\Vert\MS(\rho))=\log7-S(\rho)=\log(7/2).
\end{equation}
In particular $\rho$ is zero-mean. The largest remaining modulus in
\eqref{eq:fourier} is $\cos(\pi/7)$, so the actual factor is
$\lambda=\cos^2(\pi/7)\in(0,1)$.

For any numerical replacement $0\le\ell\le1$ of $\lambda$, put
$A=1+\tfrac52\ell$.
Then $1\le A\le7/2$ and
$B=\log7-\log A\ge\log2>0$.
At $C=\sqrt2$, the proposed right-hand side is bounded by
\[
 \frac{\log7}{B}\,\frac{\log2}{2}
 \le\frac12\log7
 <\log7-\log2.
\]
The final strict inequality is equivalent to $7>4$.
This contradicts \eqref{eq:mean} and proves the theorem.
\end{proof}

\begin{corollary}\label{cor:tensor}
For every integer $n\ge1$, the state $\rho_n=\rho^{\otimes n}$ in
$(\mathbb C^7)^{\otimes n}$ violates \eqref{eq:expected} at $C=2^{n/2}$.
\end{corollary}
\begin{proof}
Characteristic functions and convolution outputs factor under tensor products.
A product of numbers of modulus at most one has modulus one only when every
factor does. Thus $\MS(\rho_n)=I_{7^n}/7^n$, and the largest nonunit modulus
remains $\cos(\pi/7)$: attain it in one coordinate and use the origin in the
others. Hence
\[
 R_n=7^n,\quad \Tr\rho_n^2=2^{-n},\quad S(\rho_n)=n\log2,
 \quad D_n=n\log(7/2),\quad \log\delta_q[\rho_n]=\tfrac n2\log2.
\]
With $A_n=1+\lambda((7/2)^n-1)$, the denominator
$B_n=n\log7-\log A_n$ is at least $n\log2$ for every $\lambda\in[0,1]$.
The proposed bound is therefore at most $\tfrac n2\log7<D_n$.
\end{proof}

\section{The companion statements}

Remark~45 also records the weaker estimate
\begin{equation}\label{eq:weaker}
 D(\rho\Vert\MS(\rho))\ \le\
 \frac{\log R(\rho)-S(\rho)}{\log R(\rho)-\log\bigl[1+\lambda
 (R(\rho)\Tr\rho^2-1)\bigr]-S(\rho)}\log C,
 \tag{BGJ 59}
\end{equation}
under a positive-denominator hypothesis.

\begin{corollary}\label{cor:weaker}
For every $\rho_n$ in Corollary~\ref{cor:tensor}, the denominator in
\eqref{eq:weaker} is positive, and that inequality holds strictly for
every $C\ge\delta_q[\rho_n]$.
\end{corollary}
\begin{proof}
The actual factor satisfies $3/4<\lambda=\cos^2(\pi/7)<1$, since
$0<\pi/7<\pi/6$. The notation from Corollary~\ref{cor:tensor} gives
\[
 0<B_n-n\log2
 =-\log\!\left(\lambda+(1-\lambda)(2/7)^n\right)
 <-\log\lambda<\log(4/3)<\tfrac12\log2
 \le\tfrac n2\log2.
\]
Here $\log(4/3)<\tfrac12\log2$ follows by squaring from $16/9<2$.
The numerator in \eqref{eq:weaker} equals $D_n>0$ and
$\log C\ge\tfrac n2\log2$. Its right-hand side is therefore strictly
larger than $D_n$, proving the assertion without a numerical approximation.
\end{proof}

For completeness, the entropy-growth statement in the same remark holds for all
computationally diagonal states on $(\mathbb C^7)^{\otimes n}$ with $s=t=2$.
Let $X,Y$ be independent copies of any finitely
supported random variable on $\mathbb F_7^n$. Then
\[
 H(2X+2Y)\ge H(2X+2Y\mid Y)=H(X).
\]
Equality means that the conditional laws of $2X+2Y$ are identical for all
$Y$ in its support. Fix a support point $a$. After dividing by two, this says
that the law of $X$ is invariant under every translation $y-a$ with $y$ in
the support. Let $K$ be the subgroup generated by these differences.
The support lies in $a+K$, and translation invariance forces the law to be
uniform on that coset. Conversely, $2K=K$, so if $X$ is uniform on $a+K$, then $2X+2Y$
is uniform on $4a+K$ and has the same entropy. Such diagonal coset states are
minimal stabilizer-projection states. Indeed, with $K^\perp$ the annihilator of
$K$, the projection onto the coset's span is
\[
 P_{a+K}=\frac1{|K^\perp|}\sum_{p\in K^\perp}\omega^{-p\cdot a}Z^p,
 \qquad \rho_{a+K}=\frac{P_{a+K}}{|K|}.
\]
Conversely, a diagonal minimal stabilizer-projection state has only clock
coefficients in its Weyl expansion; its clock eigenvalue constraints define an
affine subspace, and its normalized projection is uniform there.
This proves the equality classification within the computationally diagonal
class at these parameters, as mentioned in Remark~45, consistently
with the classical zero-doubling characterization~\cite[Theorem~1.11(i), \S5]{Tao}.

The counterexample family specifically violates~(58), satisfies~(59), and has
strict entropy growth; the source's pure-state Theorem~44 is not refuted.
The complementary quantum-difference constant of Remark~43 also
equals $\sqrt2$ for the single-qudit example: the law of $-2I+2J$ has weights
$1/2,1/4,1/4$ at $0,2,5$, respectively. A bound with the same right-hand side
and that constant would fail for the same reason.

\paragraph{Verification.}
The proofs above cover every $n\ge1$. An exact-arithmetic checker in the
accompanying source archive verifies the seven-dimensional permutation, both
convolution marginals, the characteristic-function modulus-one set in
$\mathbb Q(e^{2\pi i/7})$, and tensor-output entropies for $1\le n\le8$.
The archive's member manifest fixes the checker version; its finite checks
corroborate, rather than establish, the all-$n$ proofs.

\begin{thebibliography}{9}
\setlength{\itemsep}{0pt}
\bibitem{BGJ} K.~Bu, W.~Gu and A.~Jaffe,
\emph{Quantum Ruzsa Divergence to Quantify Magic},
IEEE Transactions on Information Theory \textbf{71} (2025), 2726--2740.
\href{https://doi.org/10.1109/TIT.2025.3543276}{doi:10.1109/TIT.2025.3543276}.
The numbered statements used here refer to
\href{https://arxiv.org/abs/2401.14385v3}{arXiv:2401.14385v3} (23 March 2026).
\bibitem{DQG} K.~Bu, W.~Gu and A.~Jaffe,
\emph{Discrete quantum Gaussians and central limit theorem},
Advances in Mathematics \textbf{503} (2026), 111204.
\href{https://doi.org/10.1016/j.aim.2026.111204}{doi:10.1016/j.aim.2026.111204}.
\bibitem{PNAS} K.~Bu, W.~Gu and A.~Jaffe,
\emph{Quantum entropy and central limit theorem},
Proceedings of the National Academy of Sciences \textbf{120}(25) (2023), e2304589120.
\href{https://doi.org/10.1073/pnas.2304589120}{doi:10.1073/pnas.2304589120}.
\bibitem{Tao} T.~Tao,
\emph{Sumset and inverse sumset theory for Shannon entropy},
Combinatorics, Probability and Computing \textbf{19} (2010), 603--639.
\href{https://doi.org/10.1017/S0963548309990642}{doi:10.1017/S0963548309990642}.
\bibitem{Goodarzi} M.~M.~Goodarzi,
\emph{Submodularity of entropy under quantum convolution},
\href{https://arxiv.org/abs/2609.40211v1}{arXiv:2609.40211v1} (2026).
\bibitem{Source} trureturing contributors,
\emph{MixedStateInverseSumsetRefutation}, source development,
revision \texttt{57b811c}, declaration \texttt{result} (2026).
\href{https://github.com/the-omega-institute/trureturing/blob/57b811cedb64c8065c6afac0da57d0fe2c8aa5ab/D5/S3/Quantum/Magic/MixedStateInverseSumsetRefutation.lean}{Pinned source file}.
\end{thebibliography}

\smallskip\noindent\small
\textcopyright\ 2026 Haobo Ma and Wenlin Zhang.
Licensed under Creative Commons Attribution 4.0 International (CC BY 4.0).
\end{document}
